% The prerequisites of "The third pattern", taught from the ground up.
% One file, on purpose: build with  cd paper && tectonic -X compile tutorial.tex --outdir build
\documentclass[10pt]{article}
\usepackage[margin=1.1in]{geometry}
\usepackage{amsmath}
\usepackage{amssymb}
\usepackage{amsthm}
\usepackage{booktabs}
\usepackage{xcolor}
\usepackage{microtype}
\usepackage{tikz}
\usetikzlibrary{arrows.meta,calc}
\usepackage[most]{tcolorbox}
\usepackage[colorlinks=true, linkcolor=cool, citecolor=cool, urlcolor=cool]{hyperref}

\definecolor{ink}{HTML}{15181C}
\definecolor{accent}{HTML}{C81E5A}
\definecolor{cool}{HTML}{1B6CA8}
\definecolor{warm}{HTML}{D4761A}
\definecolor{soft}{HTML}{8A94A0}

\theoremstyle{plain}
\newtheorem{fact}{Fact}[section]
\theoremstyle{definition}
\newtheorem{definition}[fact]{Definition}
\newtheorem{example}[fact]{Example}

\newtcolorbox{why}{enhanced, colback=cool!4, colframe=cool!55, title={Why the paper needs this},
  fonttitle=\bfseries\small, coltitle=white, left=5pt, right=5pt, top=4pt, bottom=4pt,
  boxrule=0.5pt, arc=1.5pt, before skip=8pt, after skip=8pt, fontupper=\small}
\newtcolorbox{checkyourself}{enhanced, colback=warm!6, colframe=warm!65, title={Check yourself},
  fonttitle=\bfseries\small, coltitle=white, left=5pt, right=5pt, top=4pt, bottom=4pt,
  boxrule=0.5pt, arc=1.5pt, before skip=8pt, after skip=8pt, fontupper=\small}
\newcommand{\answer}{\par\smallskip\noindent\textit{Answer.}\ }

\newcommand{\R}{\mathbb{R}}
\newcommand{\Z}{\mathbb{Z}}
\newcommand{\tor}{\mathbb{T}}
\newcommand{\E}{\mathbb{E}}
\newcommand{\cnt}{\xi}
\newcommand{\het}{\eta}
\newcommand{\Moire}{\textsc{Moir\'e}}
\newcommand{\paper}{\emph{The third pattern}}

\title{Beating from the ground up\\[0.3em]\large The prerequisites of \paper, taught}
\author{}
\date{Tutorial companion, \today}

\begin{document}
\maketitle

\section*{How to use this}

This is the tutorial session that goes with \paper. The paper introduces
its ideas from instances and confines outside facts to eight ``Lookup''
boxes; this document is those boxes opened up and taught in order, so that
nothing in the paper has to be taken on trust. It assumes college calculus
(derivatives, integrals, Taylor series), linear algebra (vectors, matrices,
dot products, determinants) and a first course in probability (random
variables, expectation, conditioning, independence), and nothing else: not
complex analysis, not Fourier analysis, not number theory, not topology.
Everything else is taught here from those three, and where a fact is taken
on trust --- there are four, each named where it occurs --- the text says so.

Read it beside the paper, not before it. Each part opens with what the
paper needs from it and closes with short exercises whose answers are
printed; do them with a pencil first. The map:

\begin{center}\small
\begin{tabular}{@{}lll@{}}
\toprule
Paper & Its lookup or need & This tutorial \\
\midrule
\S1--2 & the count; the torus as a glued square; slow and fast & Parts~\ref{p:count}, \ref{p:torus}, \ref{p:rates} \\
\S3 & complex waves; the one integral; Fourier series; conditional expectation & Parts~\ref{p:waves}, \ref{p:fourier}, \ref{p:average} \\
\S3 & a blur as a multiplier; the curvature remainder & Parts~\ref{p:rates}, \ref{p:blur} \\
\S4 & the waves of a pulse train; front ends; hard and soft & Parts~\ref{p:fourier}, \ref{p:frontend} \\
\S5 & the price; continued fractions; lattice reduction & Parts~\ref{p:cf}, \ref{p:lattice} \\
\S6 & products of waves; the tower property; invariants and quotients & Parts~\ref{p:waves}, \ref{p:average}, \ref{p:symmetry} \\
\S7 & the implicit function theorem; the three rungs; folds; winding & Parts~\ref{p:implicit}, \ref{p:winding} \\
\S5, \S8 & a strobe as a family: sampling and aliasing & Part~\ref{p:sampling} \\
\bottomrule
\end{tabular}
\end{center}

A glossary of the paper's own words is Part~\ref{p:glossary}.

\section{Counting, and why frequency is an average}
\label{p:count}

\begin{why}
The paper's first idea is that a family of repeated members has a
\emph{count}, and that the count is a function of position that exists
whether or not the family is regular. Everything downstream is built on
counts, not frequencies.
\end{why}

Take ticks on a ruler at positions $t_1<t_2<t_3<\dots$. The \emph{count}
$\cnt(x)$ is ``how many ticks up to $x$'', with the whole-number staircase
smoothed out: $\cnt(t_n)=n$, and between $t_n$ and $t_{n+1}$ the count rises
linearly (or smoothly) from $n$ to $n+1$. For evenly spaced ticks at pitch
$T$ (the distance between neighbours), $\cnt(x)=x/T$ exactly, and the
\emph{rate} $\cnt'(x)=1/T$ is what a physicist calls the spatial frequency.
For unevenly spaced ticks $\cnt$ is still a perfectly good increasing
function; its rate varies from place to place, and ``the frequency'' is just
the average rate over a stretch. That is the whole reason the paper never
needs a frequency: a count says which member is here, a frequency says how
many there are on average, and beating depends on the former.

The \emph{fractional part} $\{\cnt(x)\}=\cnt(x)-\lfloor\cnt(x)\rfloor$, also
written $\cnt(x)\bmod1$, says where you are within the current member: $0$
on a member, $\tfrac12$ halfway to the next. Two things to notice. First,
$\{\cnt\}$ is periodic in the count with period one, so any question about
coincidences between families is a question about fractional parts.
Second, a whole-number combination of counts is a count: if $\cnt_A$ and
$\cnt_B$ are counts, $D=\cnt_A-\cnt_B$ rises by one exactly when family A
has gained one member on family B, and its fractional part depends only on
the fractional parts of $\cnt_A$ and $\cnt_B$, because
$\{a\cnt_A+b\cnt_B\}=\{a\{\cnt_A\}+b\{\cnt_B\}\}$ for whole numbers $a,b$
(a whole-number combination of whole numbers is a whole number, and drops
out). A tick of A lies exactly on a tick of B where both fractional parts
are zero; if the two families share a tick at $x=0$, that happens exactly
where $D$ is a whole number.

\begin{checkyourself}
Ten ticks to the inch against nine, sharing a tick at $x=0$. Compute
$D(x)=\cnt_A(x)-\cnt_B(x)$ and find the period of the coincidence pattern.
\answer $\cnt_A=10x$, $\cnt_B=9x$ (inches), $D=x$: the coincidences, where
$D$ is a whole number, are one inch apart. Ten members of A and nine of B
per period, as the paper's Figure~1 draws.
\end{checkyourself}

\section{The circle and the torus}
\label{p:torus}

\begin{why}
The paper's state $\Phi(x)=(\cnt_1,\dots,\cnt_K)\bmod1$ is a point on a
torus, and the ``picture'' is a function on that torus. Slides are straight
lines on it; a slide by whole numbers of members is a closed loop; the
level lines of a combined count are straight lines too, and whether a line
closes up or winds around forever is the difference between a rational and
an irrational slope.
\end{why}

\paragraph{The circle as an interval with its ends glued.}
Numbers modulo one live on the interval $[0,1)$ with $1$ identified with
$0$: walk past $1$ and you are back at $0$. That is a circle, written
$\R/\Z$ or $\tor^1$. A function on the circle is exactly a function $f(u)$
of a real variable with $f(u+1)=f(u)$: a periodic function of period one.
Nothing about roundness is needed; only the gluing. The point $\cnt(x)\bmod1$
is where you are in the current member, and as $x$ moves the point goes
round the circle once per member.

\paragraph{The torus as a square with opposite edges glued.}
Two counts give a point $(u_1,u_2)$ of the unit square with $u_1=1$ glued to
$u_1=0$ and $u_2=1$ glued to $u_2=0$: the torus $\tor^2=\R^2/\Z^2$. A
function on $\tor^2$ is a function $I(u_1,u_2)$ periodic with period one in
each variable separately. $K$ counts give $\tor^K$, the unit cube in $K$
dimensions with opposite faces glued; you never need to picture it beyond
that.

\paragraph{Lines on the torus.}
A straight line $u(t)=u_0+tw$ in the square, with $w=(w_1,w_2)$, wraps
whenever it hits an edge. If $w$ has whole-number entries the line returns
to $u_0$ at $t=1$: it is a \emph{closed loop}. If $w_1/w_2$ is rational the
line closes after $t$ has advanced enough for both coordinates to have moved
by whole numbers; if $w_1/w_2$ is irrational, it never closes and comes
arbitrarily close to every point of the torus (Kronecker's theorem;
Part~\ref{p:symmetry} proves the stronger statement that it spends equal
time in equal areas).

\paragraph{Level lines of a combined count.}
A \emph{recipe} is a vector $k=(k_1,k_2)$ of whole numbers, and the
combined count $k\cdot u=k_1u_1+k_2u_2$ is constant along the lines
perpendicular to $k$. The line in direction $w$ leaves $k\cdot u$ unchanged
exactly when $k\cdot w=0$: that is the paper's ``recipes orthogonal to the
slide'', and it is nothing but the dot product. On the torus of Figure~2 in
the paper, the slide $w=(1,1)$ runs along the level lines of $k=(1,-1)$ and
crosses the level lines of $k=(1,1)$.

\begin{checkyourself}
Which recipes $k=(k_1,k_2)$ are unchanged along the slide $w=(1,2)$? Which
three-family recipes are unchanged along $w=(1,1,1)$?
\answer $k\cdot w=k_1+2k_2=0$: $k=(2,-1)$ and its multiples. The paper calls
$(2,-1)$ the \emph{octave} recipe: it is slow when one family's rate is
about twice the other's, the ratio of a musical octave. For $(1,1,1)$: every
$k$ with $k_1+k_2+k_3=0$, the zero-sum recipes, including $(1,-1,0)$,
$(0,1,-1)$ and $(1,-2,1)$.
\end{checkyourself}

\section{Complex numbers, waves, and the one integral}
\label{p:waves}

\begin{why}
The paper's proofs expand pictures in waves $e^{2\pi i\,k\cdot u}$ and use
one integral: over a whole number of turns a wave averages to zero unless
it does not turn at all. Its cascade proof multiplies two real waves and
reads off the sum and difference of their recipes. All of this is built
from Taylor series.
\end{why}

\paragraph{Complex numbers, in five lines.}
A complex number is $z=a+bi$ with $i^2=-1$; add and multiply as for
polynomials and replace $i^2$ by $-1$. Its \emph{conjugate} is
$\bar z=a-bi$, its \emph{modulus} is $|z|=\sqrt{a^2+b^2}$, and
$z\bar z=|z|^2$, which is how one divides: $1/z=\bar z/|z|^2$. The
conjugate of a product is the product of the conjugates, and $|zw|=|z||w|$.
Points $a+bi$ are points $(a,b)$ of the plane; $|z|$ is the distance from
the origin.

\paragraph{Euler's formula, from Taylor series.}
Put $i\varphi$ into the series $e^x=\sum x^n/n!$ and sort the terms by the
powers of $i$ ($i^2=-1$, $i^3=-i$, $i^4=1$, \dots):
\[
  e^{i\varphi}=\Big(1-\frac{\varphi^2}{2!}+\frac{\varphi^4}{4!}-\dots\Big)
  +i\Big(\varphi-\frac{\varphi^3}{3!}+\frac{\varphi^5}{5!}-\dots\Big)
  =\cos\varphi+i\sin\varphi .
\]
So $e^{i\varphi}$ is the point of the unit circle at angle $\varphi$, and
$|e^{i\varphi}|=1$. We measure angles in \emph{turns}: with
$\varphi=2\pi\theta$, the point $e^{2\pi i\theta}$ goes once round the circle
as $\theta$ runs from $0$ to $1$, and $e^{2\pi i(\theta+1)}=e^{2\pi i\theta}$:
a wave is a function on the circle of Part~\ref{p:torus}. Multiplying two
waves adds their angles, $e^{2\pi i\alpha}e^{2\pi i\beta}=e^{2\pi i(\alpha+\beta)}$,
because $e^{x+y}=e^xe^y$ holds for the series with complex arguments too.
Conjugating a wave reverses it: $\overline{e^{2\pi i\theta}}=e^{-2\pi i\theta}$.
Adding and subtracting Euler's formula for $\pm\varphi$ gives the two
identities everything below runs on:
\[
  \cos\varphi=\tfrac12\big(e^{i\varphi}+e^{-i\varphi}\big),\qquad
  \sin\varphi=\tfrac1{2i}\big(e^{i\varphi}-e^{-i\varphi}\big).
\]

\paragraph{Calculus with waves.}
A complex-valued function of a real variable, $f(u)=p(u)+iq(u)$, is
differentiated and integrated part by part, so the fundamental theorem of
calculus applies to it. Differentiating $e^{2\pi imu}=\cos2\pi mu+i\sin2\pi mu$
gives $-2\pi m\sin2\pi mu+2\pi im\cos2\pi mu=2\pi im\,e^{2\pi imu}$: the wave
is its own derivative up to the constant $2\pi im$, exactly like $e^{au}$.

\begin{fact}[The one integral]
\label{f:oneintegral}
For a whole number $m$,
\[
  \int_0^1 e^{2\pi i m u}\,du=\begin{cases}1 & m=0,\\ 0 & m\ne0.\end{cases}
\]
\end{fact}

\begin{proof}
For $m=0$ the integrand is $1$. Otherwise the antiderivative is
$e^{2\pi imu}/(2\pi im)$ by the derivative just computed, so the integral is
$(e^{2\pi im}-1)/(2\pi im)$, and $e^{2\pi im}=\cos2\pi m+i\sin2\pi m=1$ for
a whole number $m$. The picture: the average position of a point that goes
round the circle $m$ times is the centre.
\end{proof}

Two consequences the paper uses constantly. On the torus, a wave with
recipe $k$ is $e^{2\pi i\,k\cdot u}$, and along a slide $u+tw$ it becomes
$e^{2\pi i\,k\cdot u}\,e^{2\pi i\,t\,(k\cdot w)}$: a wave in $t$ turning
$k\cdot w$ times per unit of $t$, which is a whole number when $w$ is a
whole-number vector; Fact~\ref{f:oneintegral} then says its average over
one slide is itself if $k\cdot w=0$ and zero otherwise. And a real wave
$\cos2\pi\theta$ carries recipe $k$ and recipe $-k$ at once, so the product
of two real waves,
\[
  \cos2\pi\alpha\,\cos2\pi\beta=\tfrac12\big(\cos2\pi(\alpha+\beta)+\cos2\pi(\alpha-\beta)\big),
\]
(multiply out the two half-sums of exponentials) carries both the
\emph{sum} and the \emph{difference} of their recipes. That product-to-sum
identity is the whole classical theory of beats: two tones multiplied (by a
nonlinearity, Part~\ref{p:frontend}) give the sum and the difference
frequency, and the difference is the beat.

\begin{checkyourself}
The paper's cascade takes the two pairwise beats of three tones, recipes
$(1,-1,0)$ and $(0,1,-1)$, and squares their sum. The cross term of the
square is their product; which recipes does it carry?
\answer The product of two real waves carries the sum $(1,0,-1)$ and the
difference $(1,-1,0)-(0,1,-1)=(1,-2,1)$. The difference is the slow one when
the two pairwise beats have nearly equal rates: the beat of beats.
\end{checkyourself}

\section{Fourier series on the circle and on the torus}
\label{p:fourier}

\begin{why}
The picture on the torus is a sum of waves, one per recipe, and the
coefficient of each wave is its strength. A stroke's or a click's profile
has coefficients $\sin(n\pi d)/n\pi$, zero when $nd$ is a whole number:
the duty null. A product of profiles has coefficients that multiply; a sum
has no cross terms; a power of a pure wave has bounded harmonics.
\end{why}

\paragraph{On the circle.}
A periodic function $g(u)$ of period one is a sum of waves,
\[
  g(u)=\sum_{n=-\infty}^{\infty}\hat g(n)\,e^{2\pi i n u},\qquad
  \hat g(n)=\int_0^1 g(u)\,e^{-2\pi i n u}\,du .
\]
This is the first fact taken on trust: Fourier's theorem, that for a
piecewise smooth $g$ the series converges to $g$ everywhere except at a
jump, where it gives the midpoint. What is not on trust is the formula for
the coefficients, which follows from Fact~\ref{f:oneintegral}: multiply the
series by $e^{-2\pi inu}$ and integrate over one period; every term but the
$n$-th averages to zero. In the language of linear algebra the waves are
\emph{orthogonal} --- $\int_0^1e^{2\pi imu}\overline{e^{2\pi inu}}\,du$ is
$1$ for $m=n$ and $0$ otherwise --- and $\hat g(n)$ is the component of $g$
along the $n$-th wave, as $v\cdot e_n$ is the component of a vector along a
unit axis. Three facts you will use without thinking: $\hat g(0)$ is the
mean of $g$; for real $g$, $\hat g(-n)=\overline{\hat g(n)}$ (conjugate the
defining integral), so the strengths $|\hat g(n)|$ come in equal pairs and
one speaks of ``the $n$-th harmonic'' for both; and the \emph{line} at
frequency $n$ in a signal's spectrum is $|\hat g(n)|$.

\begin{example}[The pulse train]
\label{e:pulse}
Let $g(u)=1$ for $|u|<d/2$ (mod $1$) and $0$ otherwise: a stroke, or a
click, covering a fraction $d$ of each period, its \emph{duty}. Then
$\hat g(0)=d$ and for $n\ne0$
\[
  \hat g(n)=\int_{-d/2}^{d/2}e^{-2\pi inu}\,du
  =\Big[\frac{e^{-2\pi inu}}{-2\pi in}\Big]_{-d/2}^{d/2}
  =\frac{e^{-\pi ind}-e^{\pi ind}}{-2\pi in}=\frac{2i\sin(n\pi d)}{2\pi in}=\frac{\sin(n\pi d)}{n\pi},
\]
using the sine identity of Part~\ref{p:waves} in the last step but one. So
$\hat g(n)=0$ exactly when $nd$ is a whole number: a $50\%$ train has no
even harmonics, a $\tfrac13$ train no third, sixth, \dots\ And
$|\hat g(n)|\le1/(|n|\pi)$: the strengths fall like $1/|n|$, slowly, which
is what a sharp edge costs (the function $\sin(\pi t)/\pi t$ is called the
\emph{sinc}). A smooth bump's coefficients fall much faster; ``anything
with edges has many harmonics'' is this comparison.
\end{example}

\paragraph{Symmetry kills harmonics.}
Write $\delta_{n0}$ for $1$ if $n=0$ and $0$ otherwise (the coefficients of
the constant function $1$). If $g(u+\tfrac12)=1-g(u)$ --- the second half
of the period is the first turned upside down, true of every $50\%$ train
however its edges are softened symmetrically --- then the shift multiplies
$\hat g(n)$ by $e^{\pi in}=(-1)^n$ and the right side has coefficients
$\delta_{n0}-\hat g(n)$, so $(-1)^n\hat g(n)=\delta_{n0}-\hat g(n)$; for
even $n\ne0$ that says $\hat g(n)=-\hat g(n)=0$. Softening the edges by an
even smoothing preserves the symmetry, so it preserves the null: the
paper's Corollary~4.3.

\paragraph{On the torus.}
A function $I(u_1,u_2)$ periodic in each variable is a double series
\[
  I(u)=\sum_{k\in\Z^2}\hat I(k)\,e^{2\pi i\,k\cdot u},\qquad
  \hat I(k)=\int_{\tor^2}I(u)\,e^{-2\pi i\,k\cdot u}\,du,
\]
one wave per recipe $k=(k_1,k_2)$, and the same in $K$ variables with
$k\in\Z^K$. Four structural facts:
\begin{itemize}
\item \emph{A product of profiles.} If $I(u_1,u_2)=g_1(u_1)g_2(u_2)$ then
$\hat I(k_1,k_2)=\hat g_1(k_1)\hat g_2(k_2)$: multiply the two series and
collect terms. Every recipe is present, with the product of the two
harmonics' strengths. This is the printed pair of gratings, whose inks
multiply.
\item \emph{A sum of profiles.} If $I=g_1(u_1)+g_2(u_2)$ then $\hat I(k)$ is
nonzero only for $k=(n,0)$ or $(0,n)$: each profile's series has waves in
its own variable alone, and a sum adds series. No cross recipe. This is two
sounds in the air, or two projected gratings.
\item \emph{A polynomial of a sum.} $(g_1+g_2)^2=g_1^2+2g_1g_2+g_2^2$: the
cross term $g_1g_2$ is a product, so the square of a sum carries the cross
recipes $(a,b)$ at strength $2\hat g_1(a)\hat g_2(b)$. In general the $n$-th
power of a sum of $K$ profiles is a sum of monomials $g_1^{e_1}\cdots g_K^{e_K}$
with $e_1+\dots+e_K=n$.
\item \emph{Powers of a pure wave.} If $g(u)=\cos2\pi u$ then
$g^e=2^{-e}(e^{2\pi iu}+e^{-2\pi iu})^e$, and the binomial expansion has
exponents $e,e-2,\dots,-e$ only: $g^e$ carries harmonics of size at most
$e$. So a monomial $g_1^{e_1}\cdots g_K^{e_K}$ of pure waves carries only
recipes with $|k_i|\le e_i$, of \emph{order} $\sum|k_i|\le n$. That bound is
the paper's Theorem~6.2, and it fails for harmonic-rich profiles, which is
why that theorem says ``pure tones''.
\end{itemize}

\paragraph{Two-valued profiles.}
If $g$ takes only the values $0$ and $1$ then $g^2=g$, $g^3=g$, and so on.
So any polynomial in such profiles collapses to a combination of the
products $g_1$, $g_2$, $g_1g_2$, $1$ --- and each of those carries at recipe
$(a,b)$ a product of the profiles' own coefficients. A harmonic that a
profile lacks cannot be manufactured by any front end: the paper's
Corollary~4.2(b), and the reason the duty null is heard through every ear.

\begin{checkyourself}
(a) At what duties does a stroke have no third harmonic? (b) Two $30\%$
pulse trains added: what is the coefficient of the recipe $(1,-1)$ in the
sum, and in the square of the sum divided by four?
\answer (a) $3d\in\Z$: $d=\tfrac13,\tfrac23$. (b) In the sum, zero. In
$(g_1+g_2)^2/4$, the cross term $2g_1g_2/4$ gives
$\tfrac12\hat g_1(1)\hat g_2(-1)=\tfrac12(\sin0.3\pi/\pi)^2=\tfrac12(0.2575)^2\approx0.033$.
\end{checkyourself}

\section{Averages along a slide, and conditional expectation}
\label{p:average}

\begin{why}
The paper's central object is $\E[\,I\mid\text{slow}\,]$: the picture
averaged over everything the slow recipes do not fix. Here is what it means
in the language of your probability course, why the sharp slide computes
it exactly, and why averages compose.
\end{why}

\paragraph{The average along a closed slide.}
Take a whole-number vector $w$ and average the picture along the closed
loop $u+tw$, $0\le t<1$:
$\mathcal E_wI(u)=\int_0^1I(u+tw)\,dt$. Expand $I$ in waves; each wave
becomes $e^{2\pi i\,k\cdot u}e^{2\pi it(k\cdot w)}$, and by
Fact~\ref{f:oneintegral} its average is itself if $k\cdot w=0$ and zero
otherwise. So
\[
  \mathcal E_wI(u)=\sum_{k\cdot w=0}\hat I(k)\,e^{2\pi i\,k\cdot u}:
\]
the slide keeps exactly the recipes orthogonal to it, with nothing left
over. That is the paper's Theorem~3.1, and you have now proved it. A slide
by a \emph{non}-whole-number vector would turn each wave a fractional
number of times, and Fact~\ref{f:oneintegral} would not apply; the average
would be a little of every wave.

\paragraph{Conditional expectation, as a probabilist says it.}
Let $U$ be a point chosen uniformly at random on the torus (every region
gets probability equal to its area). The picture's value $I(U)$ is a random
variable, and $\E[I(U)]=\hat I(0)$ is its mean. Now condition on the slow
combined counts: let $S$ be the values $k\cdot U\bmod1$ for the recipes $k$
in the kept set $L=\{k:k\cdot w=0\}$. However your course defined
$\E[X\mid S]$ for a continuous $S$ --- through conditional densities, or
as the function of $S$ that predicts $X$ best in mean square; the two
agree --- it comes to this: $\E[I(U)\mid S]$ is the average of $I$ over the
set of points $U$ that share the given values of $S$, with the density $U$
has on that set, which for a uniform $U$ is uniform along it. Which points
share the same slow counts? Take $w$ \emph{primitive} (entries with no
common factor; $(2,2)$ traces the loop of $(1,1)$ twice and says nothing
new). Then the whole-number recipes orthogonal to $w$ are the multiples of
one recipe $k_0$, the condition is $k_0\cdot u\equiv k_0\cdot u_0\pmod1$,
and that set is a family of parallel lines in the square which the
wrapping loop $u_0+tw$ visits every one of before it closes (try
$w=(2,3)$, $k_0=(3,-2)$: the loop crosses the lines $3u_1-2u_2=m$ for
$m=0,1,\dots$ as it wraps). So the points sharing the slow counts are
exactly the closed slide through $u_0$, and the conditional expectation
given $S$ is the average along the slide: $\E[I\mid S]=\mathcal E_wI$. That
is why the paper writes $\E[\,I\mid\text{slow}\,]$ and means the slide
average.

\paragraph{The projection view.}
In the space of functions on the torus with the inner product
$\langle f,g\rangle=\int f\bar g$, the waves are orthonormal (Part~\ref{p:fourier}),
and $\E[\cdot\mid S]$ is the orthogonal projection onto the span of the
waves with recipes in $L$: it keeps those components and drops the rest.
Two consequences the paper uses. Projections are idempotent: averaging
twice changes nothing. And projecting onto a subspace and then onto a
smaller subspace inside it is projecting onto the smaller one. In
probability this is the \emph{tower property}: conditioning on less
information after conditioning on more is conditioning on less,
$\E[\E[X\mid S_1]\mid S_2]=\E[X\mid S_2]$ whenever $S_2$ is a function of
$S_1$ (the paper writes it with the information sets $\mathcal G_2\subseteq\mathcal G_1$
of $S_2$ and $S_1$; ``$\mathcal G_2\subseteq\mathcal G_1$'' means exactly
``$S_2$ carries less information than $S_1$''). For slides: keeping the
recipes with $k\cdot w=0$ and then those with $k\cdot w'=0$ keeps the
intersection, the paper's Proposition~6.1 --- averages compose, and a
cascade of poolings keeps the intersection of the recipe sets.

\paragraph{Why ``linear'' matters.}
$\E[\cdot\mid S]$ is linear: it respects sums and scalings. So any linear
rule on pictures that cannot see the fast phases factors through it (the
paper's Proposition~6.3). A nonlinear rule need not: $\int I^2$, the
picture's energy, is unchanged by every slide but is not a function of the
average --- two pictures with the same average can have different energy.

\begin{checkyourself}
Two families, $w=(1,1)$, and the picture $I(u_1,u_2)=g_1(u_1)g_2(u_2)$ of two
$50\%$ strokes, each covering $|u|<\tfrac14$. What is $\mathcal E_wI$ as a
function of $D=u_1-u_2$?
\answer Directly: sliding both strokes together, the fraction of the period
during which both are on is the overlap of two windows of width $\tfrac12$
whose centres are $D$ apart, which is $\tfrac12-|D|$ for $|D|\le\tfrac12$
(and periodic): a tent, $\tfrac12$ at $D=0$, $0$ at $D=\pm\tfrac12$. By
waves: the kept recipes are $(a,-a)$ with coefficients
$\hat g_1(a)\hat g_2(-a)=(\sin(a\pi/2)/a\pi)^2$, so
$\mathcal E_wI=\tfrac14+\sum_{a\ne0}(\sin(a\pi/2)/a\pi)^2e^{2\pi iaD}$; the
series sums to the same tent (at $D=0$ it needs
$\sum_{a\ \text{odd}>0}1/a^2=\pi^2/8$, a classical sum you may take on trust
or check numerically: the first hundred terms give $0.2495$ against the
tent's $\tfrac14$).
\end{checkyourself}

\section{Local rates, gradients, slow and fast}
\label{p:rates}

\begin{why}
``Slow'' and ``fast'' are statements about rates of combined counts; the
heterodyne ratio is a ratio of rates; and the blur theorem needs the
first-order Taylor expansion of the state map with a bound on its second
derivative.
\end{why}

\paragraph{Gradients.}
For a count $\cnt(x)$ on the plane, the gradient
$\nabla\cnt=(\partial\cnt/\partial x_1,\partial\cnt/\partial x_2)$ points in
the direction of fastest increase, and its length is the rate of increase
per unit distance in that direction: members per unit length, across the
members. It is perpendicular to the level lines of $\cnt$, which are the
members themselves. For parallel lines of pitch $T$ turned to angle
$\alpha$, $\nabla\cnt=(\cos\alpha,\sin\alpha)/T$; for rings of pitch $s$
about $c$, $\cnt=|x-c|/s$ and $\nabla\cnt=(x-c)/(s|x-c|)$, the unit vector
from the centre over $s$ (differentiate $\sqrt{(x_1-c_1)^2+(x_2-c_2)^2}$).
Rates add: the combined count $k\cdot\cnt=\sum k_i\cnt_i$ has gradient
$\sum k_i\nabla\cnt_i$. The $a$-th harmonic of family $1$, the wave
$e^{2\pi i\,a\cnt_1}$, repeats $a$ times per member, so its rate is
$a\nabla\cnt_1$.

\begin{definition}[The heterodyne ratio, and the beat regime]
For a recipe $(a,b)$ on two families,
\[
  \het_{a,b}=\frac{\lvert a\nabla\cnt_1+b\nabla\cnt_2\rvert}
                  {\tfrac12\big(\lvert a\nabla\cnt_1\rvert+\lvert b\nabla\cnt_2\rvert\big)},
\]
the rate of the combined count over the average rate of the two harmonics
that carry it (the $\tfrac12$ makes the denominator an average of two
things). Its reciprocal is the number of carrier members per band. The
paper calls $\het<\tfrac14$ the \emph{beat regime}: more than four members
per band. The number four is a convention, not a theorem --- it is where
the paper draws the line between a band, wider than the members that make
it, and a texture; the paper's ratio view lets you see the line drawn.
\end{definition}

The word is the radio engineer's: to \emph{heterodyne} is to multiply two
signals and keep the difference frequency (Part~\ref{p:waves}'s
product-to-sum identity), and the ratio says how much slower the
difference is than its ingredients.

\begin{example}
Two families of parallel lines, pitch $1$, one turned by $10^\circ$, so
that $\nabla\cnt_1=(1,0)$ and
$\nabla\cnt_2=(\cos10^\circ,\sin10^\circ)=(0.985,0.174)$. The difference
recipe $(1,-1)$ has rate $|(0.015,-0.174)|=0.174$; the average member rate
is $1$; $\het=0.174$, about $5.7$ members per band, inside the beat regime.
Turn the second family to $20^\circ$ and $\het=0.347$: fewer than four
members per band, and the bands stop reading as bands.
\end{example}

\paragraph{The state map near a point.}
Write $\Phi=(\cnt_1,\dots,\cnt_K)$ and let $J$ be the $K\times2$ matrix
whose rows are the gradients $\nabla\cnt_i(x)$ (the \emph{Jacobian}).
Taylor's theorem in two variables says
\[
  \Phi(y)=\Phi(x)+J\,(y-x)+E(y),\qquad |E(y)|\le\tfrac12\kappa\,|y-x|^2,
\]
where $\kappa$ is the largest absolute value that any second partial
derivative of any count takes near $x$. So a small disc around $x$ maps, to
first order, onto a small parallelogram around $\Phi(x)$ spanned by the
columns of $J$ --- and when all the families have nearly the same pitch and
direction, the rows of $J$ are nearly equal and the parallelogram collapses
to a short segment in the direction $(1,1,\dots,1)$: the black segment of
the paper's Figure~2, the ``fast direction''. A combined count $k\cdot\Phi$
changes across the disc at rate $J^{\top}k=\sum k_i\nabla\cnt_i$, which is
just the rate of the combination again: $k\cdot J(y-x)=(J^{\top}k)\cdot(y-x)$.
The error term $E$ is what curvature of the members costs, and it is
second order in the size of the disc.

\begin{checkyourself}
Two ring families of pitch $s$ centred at $c_1$ and $c_2$. Where is the
difference recipe slow?
\answer Its rate is $|\nabla\cnt_1-\nabla\cnt_2|/s$, the distance between
two unit vectors pointing away from the two centres. It is exactly zero on
the line through the centres, outside the segment between them, where the
two unit vectors agree; small far from both centres in any direction; and
large between the centres, where the vectors oppose.
\end{checkyourself}

\section{Blurs, convolution, and the response of a window}
\label{p:blur}

\begin{why}
A pooling observer averages over a neighbourhood with a weight $W$. The
paper's Theorem~3.3 says this multiplies each recipe's wave by the window's
response at that recipe's local rate. Here is what a convolution is, what
the response is, why a Gaussian blur is a low-pass filter, and where the
remainder bound comes from.
\end{why}

\paragraph{Convolution.}
Blurring a signal $S(x)$ on the plane with a weight $W$ (nonnegative, total
weight one, symmetric about zero) means replacing each value by the
weighted average of its neighbours:
\[
  (S*W)(x)=\int S(y)\,W(x-y)\,dy=\int S(x-z)\,W(z)\,dz .
\]
The second form says: look a displacement $z$ away, weight by $W(z)$, add
up. Averaging in time with a slow meter is the same formula in $t$. A
window of width $\rho$ is the unit window stretched, $W_\rho(z)=\rho^{-2}W(z/\rho)$
in the plane (the factor keeps the total weight one).

\paragraph{The response.}
Feed a plain wave $S(x)=e^{2\pi i\nu\cdot x}$, with a rate vector $\nu$
that need not be a whole number, into the blur:
\[
  (S*W)(x)=\int e^{2\pi i\nu\cdot(x-z)}W(z)\,dz=e^{2\pi i\nu\cdot x}\int e^{-2\pi i\nu\cdot z}W(z)\,dz
  =\widehat W(\nu)\,e^{2\pi i\nu\cdot x}.
\]
The wave comes out as itself times a number $\widehat W(\nu)$, the window's
\emph{response} at rate $\nu$ (real, by the symmetry of $W$). This
$\widehat W$ is the \emph{Fourier transform} of $W$: the same hat as a
Fourier coefficient, but for a function on the whole plane rather than a
periodic one, it is a function of a continuous rate rather than a list
indexed by whole numbers. The two are cousins, not the same object.
Substituting $z=\rho z'$ in the integral for the stretched window gives
$\widehat{W_\rho}(\nu)=\widehat W(\rho\nu)$, the form the paper writes. Two
properties: $\widehat W(0)=\int W=1$, a constant passes unchanged; and for
$|\nu|$ large compared with $1/\rho$ the integral averages a fast wave
against a slowly varying weight and is nearly zero. In between, the
response falls off on the scale $|\nu|\sim1/\rho$: a \emph{low-pass}
filter, passing rates below about $1/\rho$ and blocking those above.

\begin{example}[The Gaussian]
On the line, $W(z)=e^{-z^2/2}/\sqrt{2\pi}$ has response
$\widehat W(\nu)=e^{-2\pi^2\nu^2}$. Complete the square in the exponent,
$-z^2/2-2\pi i\nu z=-(z+2\pi i\nu)^2/2-2\pi^2\nu^2$, so that
\[
  \widehat W(\nu)=e^{-2\pi^2\nu^2}\int e^{-(z+2\pi i\nu)^2/2}\,\frac{dz}{\sqrt{2\pi}},
\]
and the remaining integral is $\sqrt{2\pi}/\sqrt{2\pi}=1$: shifting a Gaussian by a
purely imaginary amount does not change its integral, the second fact
taken on trust (it is a short contour argument in complex analysis, and you
can check it numerically). In the plane the Gaussian is a product of two
line Gaussians and so is its response: $e^{-2\pi^2|\nu|^2}$, and the
stretched window has $e^{-2\pi^2\rho^2|\nu|^2}$. A flat window, one that is
constant on an interval, has a sinc response, Example~\ref{e:pulse} again.
\end{example}

\paragraph{The second moment.}
$m_2=\int|z|^2W(z)\,dz$ measures the unit window's width squared (for the
plane Gaussian above, $m_2=2$), and the stretched window's second moment
is $m_2\rho^2$. It enters the paper's remainder bound because the
curvature error $E(y)$ of Part~\ref{p:rates} is at most $\tfrac12\kappa|y-x|^2$,
and the average of $|y-x|^2$ under $W_\rho$ is exactly $m_2\rho^2$.

\paragraph{The blur theorem, assembled.}
Expand $S=I\circ\Phi$ in waves: $S(y)=\sum_k\hat I(k)e^{2\pi i\,k\cdot\Phi(y)}$.
Near $x$, $k\cdot\Phi(y)=k\cdot\Phi(x)+\nu_k\cdot(y-x)+k\cdot E(y)$ with
$\nu_k=J^{\top}k$ the recipe's local rate. Ignoring $E$, each term is a plain
wave at rate $\nu_k$ times a constant, and the blur multiplies it by
$\widehat W(\rho\nu_k)$. Restoring $E$, the factor $e^{2\pi i\,k\cdot E}$
differs from $1$ by at most $2\pi|k||E|$ (a point on the unit circle at
angle $\varphi$ is within $|\varphi|$ of $1$), whose average under the window
is at most $2\pi|k|\cdot\tfrac12\kappa m_2\rho^2=\pi|k|\kappa m_2\rho^2$;
summing over recipes gives the paper's remainder bound. In the beat regime
the slow recipes have $\rho|\nu_k|\ll1$ and are passed, the fast ones have
$\rho|\nu_k|\gg1$ and are blocked, and the blur reports the average of
Part~\ref{p:average} with a little leakage from the window's shoulders.

\begin{checkyourself}
A Gaussian blur of width $\rho=2$ pixels on a family of pitch $3$ pixels
and its difference beat of pitch $30$ pixels. What fraction of each survives?
\answer Rates $\tfrac13$ and $\tfrac1{30}$ per pixel;
$e^{-2\pi^2\cdot4\cdot\nu^2}$ gives $e^{-8.77}\approx1.5\times10^{-4}$ for
the members and $e^{-0.088}\approx0.92$ for the beat.
\end{checkyourself}

\section{Front ends, composition, and hard versus soft}
\label{p:frontend}

\begin{why}
An observer applies a function $N$ to each value before pooling. The paper's
one-line lemma $N\circ(I\circ\Phi)=(N\circ I)\circ\Phi$ is associativity of
composition; what $N$ can and cannot do to the recipes is
Part~\ref{p:fourier}'s polynomial expansion.
\end{why}

Composition of functions is associative: $(N\circ I)\circ\Phi$ and
$N\circ(I\circ\Phi)$ both mean ``first $\Phi$, then $I$, then $N$''. So a
front end applied to the superposition is the same as a front end applied
to the picture, looked up at the same state. The state map, and every rate,
is untouched.

What a front end does to recipes is Part~\ref{p:fourier}: a polynomial $N$
of degree $n$ applied to an additive picture produces monomials of total
degree at most $n$, hence (for pure tones) recipes of order at most $n$;
applied to two-valued profiles it produces nothing a product of the
profiles does not already contain; applied to a multiplicative picture it
changes strengths only. Two front ends that are not polynomials, on the
\emph{mean} picture $(g_1+g_2)/2$ of two two-valued profiles, which takes
the values $0$, $\tfrac12$, $1$: a saturation $\min(g_1+g_2,1)$ equals
$g_1+g_2-g_1g_2$, the picture of two inks laid one over the other (paint
covers paint), which is why a saturating eye sees an added pair as if it
were printed; and a threshold $N(v)=[v>\tfrac12]$ fires only where the mean
is $1$, both profiles on: the coincidence detector, whose picture is the
product $g_1g_2$.

\paragraph{Pool then respond, or respond then pool.}
Applying a function $N$ \emph{after} pooling gives $N(\E[I\mid S])$, a
function of the slow counts, so it carries only slow recipes, though in
new proportions ($N(e^{2\pi iD})$ contains $e^{2\pi i\,2D}$, recipe
$(2,-2)$, still slow). Applying $N$ \emph{before} pooling gives
$\E[N\circ I\mid S]$, which can carry recipes that $I$ never had --- an
additive $I$ has no cross recipe, $I^2$ does. Both are averages of a
picture along the same fast directions.

\begin{checkyourself}
Which of these front ends, applied to the mean of two added hard $50\%$
pulse trains at a mistuned octave (one train's rate about twice the
other's), lets a low-pass hear the $(2,-1)$ beat: identity, square, cube,
threshold? And after the edges are softened symmetrically?
\answer Hard: none. The beat needs the lower train's second harmonic, which
a $50\%$ train lacks, and two-valued profiles cannot be given a harmonic by
any front end (Part~\ref{p:fourier}). Softened: the square still cannot
(its cross term $g_1g_2$ multiplies one harmonic of each train, and
$\hat g_1(2)$ is still zero by the symmetry argument), but the cube can:
writing $g_1=\tfrac12+h$ with $h$ carrying only odd harmonics, $h^2$ has
even ones, so the cube's term $g_1^2g_2$ carries $(2,-1)$; for the hard
train $h^2=\tfrac14$ is constant and nothing appears.
\end{checkyourself}

\section{Continued fractions and best approximation}
\label{p:cf}

\begin{why}
Which recipe is slow, for two families at pitch ratio $x$, is a question
about how well $x$ is approximated by fractions. The paper prices a recipe
by the product of its entries, calls a high-order recipe that shows a
\emph{station}, and uses the error formula for convergents and Hurwitz's
constant.
\end{why}

\paragraph{The price, and stations.}
A recipe $(a,-b)$ between two families is carried by the $a$-th harmonic of
one and the $b$-th of the other (Part~\ref{p:fourier}, products), each of
strength about $1/n$ for a family with edges; so its beat has about
$1/|ab|$ of the contrast of the first-order beat, and to look as strong it
must be $|ab|$ times slower. The paper's \emph{priced merit} is
$\mu_{a,b}=\het_{a,b}\cdot|ab|$, the visible recipe is the one with the
smallest $\mu$, and a recipe beyond first order that shows, $\mu<\tfrac14$,
is a \emph{station}: it locks the two families at a rational ratio of
pitches. Without the price every real ratio has fractions as close as you
like, and some recipe of enormous order would always be slowest.

\paragraph{The expansion.}
For a real $x>0$, write $x=a_0+1/x_1$ with $a_0=\lfloor x\rfloor$ and
$x_1>1$; then $x_1=a_1+1/x_2$, and so on. The whole numbers $a_0;a_1,a_2,\dots$
are the \emph{partial quotients} and $x_n$ the \emph{complete quotients},
$x_n=[a_n;a_{n+1},\dots]$. Rational $x$ terminates; irrational $x$ does not.
Examples: $\pi=[3;7,15,1,292,\dots]$; $\sqrt2=[1;2,2,2,\dots]$; the silver
ratio $1+\sqrt2=[2;2,2,2,\dots]$; the golden ratio
$\varphi=[1;1,1,1,\dots]$.

\paragraph{Convergents.}
Truncating gives fractions $h_n/k_n=[a_0;a_1,\dots,a_n]$, computed by the
recurrences
\[
  h_n=a_nh_{n-1}+h_{n-2},\qquad k_n=a_nk_{n-1}+k_{n-2},\qquad
  (h_{-1},k_{-1})=(1,0),\ (h_{-2},k_{-2})=(0,1).
\]
For $\pi$: $3/1,\ 22/7,\ 333/106,\ 355/113,\dots$ Two identities, both by
induction on the recurrences: $h_{n-1}k_n-h_nk_{n-1}=(-1)^n$, and
$x=(x_{n+1}h_n+h_{n-1})/(x_{n+1}k_n+k_{n-1})$ (replace the tail of the
expansion by its exact value $x_{n+1}$; here $k_n$ is a denominator, not a
recipe).

\begin{fact}[The error of a convergent]
\label{f:cferror}
$\displaystyle |k_nx-h_n|=\frac{1}{x_{n+1}k_n+k_{n-1}}$.
\end{fact}

\begin{proof}
From the second identity,
\[
  k_nx-h_n=\frac{k_n(x_{n+1}h_n+h_{n-1})-h_n(x_{n+1}k_n+k_{n-1})}{x_{n+1}k_n+k_{n-1}}
  =\frac{h_{n-1}k_n-h_nk_{n-1}}{x_{n+1}k_n+k_{n-1}}=\frac{(-1)^n}{x_{n+1}k_n+k_{n-1}}.\qedhere
\]
\end{proof}

So a convergent is a good approximation exactly when the \emph{next}
complete quotient is large: $22/7$ is excellent for $\pi$ because
$x_2=15.996$, and $355/113$ is famous because $x_4=292.6$. Two classical
theorems, the third and fourth facts taken on trust: every convergent is a
\emph{best approximation}, no fraction with a smaller denominator is closer
(Lagrange); and for every irrational $x$ infinitely many fractions satisfy
$|x-h/k|<1/(\sqrt5k^2)$, while no constant larger than $\sqrt5$ works for
every $x$ (Hurwitz). The number that makes $\sqrt5$ sharp is the golden
ratio, whose partial quotients are all $1$, the smallest possible: it is
the \emph{worst approximable} number, and Fact~\ref{f:cferror} shows why,
because $k_n|k_nx-h_n|=1/(x_{n+1}+k_{n-1}/k_n)\to1/(\varphi+1/\varphi)=1/\sqrt5$
for it. Numbers whose partial quotients stay bounded are \emph{badly
approximable}; $\sqrt2$ and the silver ratio are examples.

\paragraph{From approximation to slowness.}
Two parallel families of pitches $s_1$ and $s_2$, ratio $x=s_1/s_2$ (coarse
over fine, so $x>1$). The recipe $(h,-k)$ has rate $h/s_1-k/s_2=(h-kx)/s_1$,
so it is slow exactly when $h/k$ approximates $x$ well, and the recipes that
beat every lower-order recipe are the convergents. With
Fact~\ref{f:cferror} the priced merit of $(h_n,-k_n)$ is
\[
  \mu_n=\frac{|h_n/s_1-k_n/s_2|}{\tfrac12(h_n/s_1+k_n/s_2)}\,h_nk_n
  =\frac{2h_nk_n\,|h_n-k_nx|}{h_n+k_nx}
  =\frac{2h_nk_n}{(h_n+k_nx)(x_{n+1}k_n+k_{n-1})},
\]
which is the paper's closed form. Since $h_n\approx k_nx$, this is about
$k_n|k_nx-h_n|=1/(x_{n+1}+k_{n-1}/k_n)$: a station exists, $\mu_n<\tfrac14$,
when the next complete quotient (plus a fraction below one) exceeds about
$4$. For the golden ratio $\mu_n\to1/\sqrt5\approx0.447$: never a station,
Hurwitz's constant reappearing as a desert.

\begin{checkyourself}
Find the convergents of $2.05=41/20$ and say which is a station.
\answer $2.05=[2;20]$: two convergents, $2/1$ and $41/20$, and the expansion
ends. For $2/1$ the next complete quotient is $20$, so $\mu_0\approx1/20$: a
strong station, the octave. The second convergent is the ratio itself, with
rate exactly zero --- not a slow beat but a constant, which no slide can
remove; the paper's tool gives such an exact lock a small floor merit so it
cannot out-pick a visible band. The visible station is the octave, and the
$16.4{:}8$ pair of the paper is this ratio.
\end{checkyourself}

\section{Lattices in the plane, and Lagrange--Gauss reduction}
\label{p:lattice}

\begin{why}
In the plane the rates of the two families are vectors, the combined
counts' rates $a\nabla\cnt_1+b\nabla\cnt_2$ form a lattice, and the slow
recipe is a short vector of that lattice. The paper finds it by
``Lagrange--Gauss reduction, the two-dimensional Euclidean algorithm''.
\end{why}

\paragraph{Lattices.}
Two vectors $v_1,v_2$ in the plane, not parallel, generate the lattice
$\Lambda=\{av_1+bv_2:a,b\in\Z\}$, a regular grid of points. Many pairs
generate the same grid: $(v_1,v_2)$ and $(v_1,v_2+v_1)$ do, and in general
any pair obtained by whole-number combinations with determinant $\pm1$
(a \emph{unimodular} change of basis). The grid has a well-defined
\emph{shortest nonzero vector}, and a basis in which $v_1$ is a shortest
vector and $v_2$ is as short as it can then be is called \emph{reduced}.
The area of the basic parallelogram, $|\det(v_1,v_2)|$, is the same for
every basis.

\paragraph{The reduction algorithm.}
Given $(u,v)$ with $|u|\le|v|$: replace $v$ by $v-mu$ with
$m=\operatorname{round}(u\cdot v/|u|^2)$, the whole number that makes the
new $v$ as short as possible in the direction of $u$; if the new $v$ is
shorter than $u$, swap them; repeat until nothing changes. It is the
Euclidean algorithm with vectors: subtract the right multiple, swap,
repeat. It stops when $|u|\le|v|$ and $|u\cdot v|\le\tfrac12|u|^2$.

\begin{fact}[A reduced basis begins with a shortest vector]
If $|u|\le|v|$ and $|u\cdot v|\le\tfrac12|u|^2$, then every nonzero lattice
vector $au+bv$ has length at least $|u|$, and at least $|v|$ if $b\ne0$.
\end{fact}

\begin{proof}
Expand and use $|u\cdot v|\le\tfrac12|u|^2$:
\[
  |au+bv|^2=a^2|u|^2+2ab\,u\cdot v+b^2|v|^2\;\ge\;a^2|u|^2-|ab||u|^2+b^2|v|^2 .
\]
If $|a|\ge|b|$ the first two terms are at least $0$ and the whole is at least
$b^2|v|^2$, which is at least $|v|^2\ge|u|^2$ when $b\ne0$, and at least
$a^2|u|^2\ge|u|^2$ when $b=0$. If $|b|>|a|$, use $|u|\le|v|$ on the middle
term: the whole is at least $|v|^2(b^2-|ab|)+a^2|u|^2\ge|v|^2|b|(|b|-|a|)\ge|v|^2$.
\end{proof}

In one dimension, with $u=1/s_1$ and $v=1/s_2$ for $s_1>s_2$ (so $|u|\le|v|$),
the rounds \emph{are} the continued fraction of $s_1/s_2$ and the
intermediate vectors are the rates of the convergents' recipes.

\begin{example}
The families of Part~\ref{p:rates}: $u=\nabla\cnt_1=(1,0)$,
$v=\nabla\cnt_2=(0.985,0.174)$. Then $m=\operatorname{round}(0.985)=1$,
$v\leftarrow v-u=(-0.015,0.174)$, shorter than $u$, swap:
$u=(-0.015,0.174)$, $v=(1,0)$. Now $u\cdot v/|u|^2=-0.015/0.0305=-0.49$,
$m=0$, nothing changes: done. The shortest vector is
$\nabla\cnt_2-\nabla\cnt_1$, the recipe $(-1,1)$, the difference beat, at
rate $0.174$: what the eye sees when a sheet is turned ten degrees, found
now by arithmetic rather than by the heterodyne ratio.
\end{example}

The paper's shader does this at every pixel, then checks a small window of
combinations around the reduced basis with the price $|ab|$ attached,
because the cheapest \emph{priced} recipe need not be the shortest vector
(a longer vector with smaller whole-number coefficients can win). The
one-dimensional statement, that the winners are convergents, is the
special case in which all the rate vectors are parallel.

\begin{checkyourself}
Reduce the lattice generated by $u=(1,0)$ and $v=(0.5,0.02)$. What recipe
wins, and what is its rate?
\answer $m=\operatorname{round}(0.5)$ is a tie; take $1$:
$v\leftarrow(-0.5,0.02)$, shorter, swap; then $u\cdot v/|u|^2=-0.5/0.2504\approx-2.0$,
$v\leftarrow v+2u=(0,0.04)$, shorter, swap: $u=(0,0.04)$, $v=(-0.5,0.02)$;
$u\cdot v/|u|^2=0.0008/0.0016=0.5$, another tie, and either choice leaves
$|v|\approx0.5$: done (the other branch at the first tie reaches the same
shortest vector). It is $2v_{\text{orig}}-u_{\text{orig}}$, the recipe
$(-1,2)$: the second family's rate is about half the first's, so its pitch
is about twice the first's, and the octave beat, at rate $0.04$, is what
shows.
\end{checkyourself}

\section{Implicit functions, the three rungs, and folds}
\label{p:implicit}

\begin{why}
A count is the solution $c$ of ``$x$ is on member $c$'', $F(x,c)=0$. Whether
that solution is a smooth function of $x$ is the implicit function theorem;
the three ways it can exist are the paper's three \emph{rungs}; where it
fails, the family folds, and the fold law is a computation with a shape's
own ruler.
\end{why}

\paragraph{Level sets and gradients.}
The set $\{x:f(x)=v\}$ for a smooth $f$ on the plane is, wherever
$\nabla f\ne0$, a smooth curve, and $\nabla f$ is perpendicular to it. The
members of a family are the level sets of its count at whole-number
values; the bands of the paper are level sets of a combined count.

\begin{fact}[Implicit function theorem]
Let $F(x,c)$ be smooth in $x\in\R^2$ and $c\in\R$, with $F(x_0,c_0)=0$ and
$\partial F/\partial c\,(x_0,c_0)\ne0$. Then near $x_0$ there is a unique
smooth function $c(x)$ with $c(x_0)=c_0$ and $F(x,c(x))=0$, and its
gradient is $\nabla c=-\nabla_xF/(\partial F/\partial c)$.
\end{fact}

The partial derivative $\partial F/\partial c$ is the rate at which $F$
changes when $c$ is nudged with $x$ held still. The idea of the proof: near
the point, $F(x,c)\approx\nabla_xF\cdot(x-x_0)+(\partial F/\partial c)(c-c_0)$,
and if the last coefficient is nonzero this linear equation can be solved
for $c$; the theorem says the nonlinear one can too, nearby. The gradient
formula is that linear solve.

\begin{example}
Ruler: $F(x,c)=x/T-c$, $\partial F/\partial c=-1\ne0$ everywhere,
$c(x)=x/T$. Rings: $F(x,c)=|x|-cs$, $c=|x|/s$, fine away from the centre.
Rays: $F(x,c)=\operatorname{angle}(x)-c\,\alpha$ with $\alpha$ the step
between rays; $c$ is the polar angle over $\alpha$, fine \emph{locally}
everywhere except the centre, but the polar angle is not a single function
on the plane: Part~\ref{p:winding}.
\end{example}

\paragraph{The three rungs.}
The paper's trichotomy says a count can exist in exactly three ways, and
now you can see why. Where $\partial F/\partial c\ne0$ everywhere, the
local solutions either glue into one function on the whole plane (the
\emph{exact} rung: rulers, rings, spirals, waves) or glue only up to a
whole number gained around a loop, because the local pieces return to
themselves shifted by an index (the \emph{winding} rung: rays, and any
family built on the polar angle; Part~\ref{p:winding}). Where
$\partial F/\partial c=0$ somewhere, the theorem fails there and the count
branches (the \emph{fold} rung, below). A field --- any function $\phi(x)$
added to the count, replacing $F(x,c)$ by $F(x,c-\phi(x))$ --- leaves
$\partial F/\partial c$ unchanged, so it can move a family between the
first two rungs and never onto the third.

\paragraph{The fold.}
Where $\partial F/\partial c=0$ the typical failure is a double root. The
model is $F(x,c)=c^2-x$ in one variable: for $x>0$ there are two solutions
$c=\pm\sqrt x$, for $x<0$ none, and at $x=0$ they meet with vertical
tangent. The solution set is a smooth curve (the parabola $x=c^2$) that is
not the graph of a function of $x$: it \emph{folds} over. In a family of
curves the fold is the \emph{envelope}: the curve that every member is
tangent to, where consecutive members meet. A nest of rings, each turned
a little more than the last, is clean near the centre and folds where a
corner of one ring first crosses the next; the paper computes that radius.

\paragraph{Fronts, gauges, and the Mach condition.}
A convex shape $K$ with the origin inside defines a \emph{gauge}
$\gamma(v)=\min\{t>0: v\in tK\}$, the scale of the copy of $K$ that just
reaches $v$: the shape's own ruler (for the disc, ordinary length; for a
square of half-side one, $\max(|v_1|,|v_2|)$). A \emph{front} is a family
$\{\gamma(x-c\delta)=cs\}$: copies of $K$ of size $cs$ centred at $c\delta$,
member $c$ being the wave that left a source at $c\delta$ and has spread
to size $cs$. Here $\partial F/\partial c$ for $F(x,c)=\gamma(x-c\delta)-cs$
is $-\nabla\gamma(x-c\delta)\cdot\delta-s$, and it vanishes somewhere on a
member exactly when $\nabla\gamma\cdot\delta=-s$ can be reached there. For
the disc, $\nabla\gamma$ is the outward unit normal $n$ of the ring, and
$n\cdot\delta$ ranges over $[-|\delta|,|\delta|]$ as $n$ goes round: so a
ring family folds if and only if $|\delta|>s$ --- the source outruns its
own waves, the \emph{supersonic} case, and the fold is the shock along the
wake, the Mach cone of a jet. Subsonic sources nest forever. For a general
shape the same computation runs with $\gamma$ in place of length and the
condition becomes $\gamma(-\delta)>s$; the companion paper does that
computation, and the disc case above is the whole idea.

\begin{checkyourself}
Rings of pitch $s=1$ whose centres step by $\delta=(1.2,0)$ per member. Do
they fold? What if $\delta=(0.8,0)$?
\answer $|\delta|=1.2>1$: supersonic, they fold along straight wake lines
from the first member. At $0.8$: subsonic, they nest.
\end{checkyourself}

\section{Winding and monodromy}
\label{p:winding}

\begin{why}
The polar angle is defined only up to whole turns; a family built on it has
a count that gains a whole number around a loop, a \emph{defect} with a
\emph{charge}, and the bands of a pair containing it must end. The paper's
winding rung is this part.
\end{why}

\paragraph{The polar angle is not a function.}
Around the origin of the plane the angle $\vartheta(x)$ can be chosen
continuously along any path, but going once round the origin brings it back
to its start \emph{plus one full turn}. No continuous choice on the whole
punctured plane exists: try, and somewhere a cut appears where the value
jumps by a turn. The \emph{winding number} of a loop about the origin is
the number of turns the angle gains along it; it is a whole number, it does
not change when the loop is deformed without crossing the origin (deform a
little and a whole number cannot jump), and it adds: a loop round two
points winds the sum of the windings round each.

\paragraph{The helix picture.}
The map $t\mapsto e^{2\pi it}$ wraps the line round the circle infinitely
often, like a helix over a circle. The angle ``should'' be a function on
the circle and is instead a function on the helix; going once round the
circle is moving one turn up the helix. A count built on the polar angle
lives on such a helix over the punctured plane and comes down to the plane
only up to whole numbers; the whole number a loop gains is its
\emph{monodromy} (the names ``covering space'' and ``universal cover''
belong to this picture, and nothing below needs more than the picture).

\paragraph{Why bands must end.}
Suppose the difference count $D$ of a pair is defined on the plane up to
whole numbers, and gains $q$ around a loop $\gamma$. Follow $D$ along
$\gamma$: as a real-valued lift it ends $q$ higher than it started, so it
crosses each whole-number level $q$ more times upward than downward. A
band --- a level line of $D$ at a whole number --- that crosses the loop
inward must, if it never ends, cross outward again, pairing an upward
crossing with a downward one. The $q$ unpaired crossings belong to bands
that enter and never leave: $q$ bands end inside $\gamma$, counted with
sign. Adding a field $q\vartheta/2\pi$ to one family's count is the
paper's way of making such a defect on purpose, and the number of endings
inside any loop is then $q$ times the enclosed points, exactly.

\paragraph{The core.}
Near the defect the added count changes at rate $q/2\pi r$ (the angle
changes by a full turn over a circle of circumference $2\pi r$), so the
beat between the modified family and its twin has heterodyne ratio
$\het=(q/2\pi r)/(1/s)=qs/2\pi r$, which reaches $\tfrac14$ at $r=2qs/\pi$.
Inside, the beat is not slow enough to be a band; the bands end at the
defect but are only bands outside the core.

\begin{checkyourself}
A field of $3\vartheta/2\pi$ is added to one of two identical line families
of pitch $s=10$. How many bands end inside a loop round the centre, and how
large is the core?
\answer Three (with the same sign). $r^\star=2\cdot3\cdot10/\pi\approx19$.
\end{checkyourself}

\section{Symmetry, invariants, quotients, and evenly spread lines}
\label{p:symmetry}

\begin{why}
The paper's final proposition says the third pattern is the quotient of the
picture by the fast rephasing: everything a linear observer can measure
without the fast phases is a function of the average. It also remarks that
a slide with incommensurable rates, averaged for a long time, keeps only
the constant. Both are about symmetry.
\end{why}

\paragraph{Shifts as a group; orbits; quotients.}
Sliding the picture, $I\mapsto I(\cdot+h)$, by every $h$ on a line through
the origin of the torus (whole and fractional amounts alike) is a
\emph{group} of transformations: two slides compose to a slide, every slide
has an inverse. Call the set of these slides $H$. The \emph{orbit} of a
picture is the set of all its slid copies. A picture is \emph{invariant}
under $H$ if every slide leaves it unchanged; a rule on pictures is
invariant if it is constant on each orbit. The \emph{quotient} of the
pictures by $H$ is the set of orbits: two pictures count as the same point
of the quotient exactly when one is a slid copy of the other, and a
function on the quotient is nothing but an invariant rule. That is all
``quotient by a symmetry'' means.

\paragraph{Averaging makes invariants.}
The average $\E_HI$ of a picture over all its slid copies is invariant:
sliding it again only relabels the copies. And it is the most economical
invariant: any \emph{linear} rule $F$ that is itself invariant satisfies
$F(I)=\text{average over }h\text{ of }F(I(\cdot+h))=F(\E_HI)$, because a
linear rule can be moved through an average. So invariant linear rules see
only the average; the average retains exactly what every slid copy has in
common and nothing else, which is why the paper calls it the quotient. In
wave language, $\E_H$ kills every wave the slides can turn and keeps every
wave they cannot: the projection of Part~\ref{p:average}. (The general
theorem behind this, that time averages along a motion converge to space
averages over what the motion fills, is the ergodic theorem; the case the
paper needs is proved next.)

\paragraph{Evenly spread lines: Weyl's theorem.}
A line $u_0+tw$ on the torus with $w_1/w_2$ irrational never closes. It does
more: it spends, in the long run, equal time in equal areas. Proof by waves:
for a recipe $k\ne0$, the average of $e^{2\pi i\,k\cdot(u_0+tw)}$ over
$0\le t\le T$ is $e^{2\pi i\,k\cdot u_0}$ times
$(e^{2\pi iT(k\cdot w)}-1)/(2\pi iT(k\cdot w))$, which has modulus at most
$2/(2\pi T|k\cdot w|)$ and tends to zero as $T$ grows as long as
$k\cdot w\ne0$ --- and $k\cdot w=k_1w_1+k_2w_2$ is never zero for whole
numbers $k\ne0$ when $w_1/w_2$ is irrational. So the time average of every
nonconstant wave tends to zero, hence the time average of any picture tends
to its mean $\hat I(0)$: the paper's ``a slide whose rates have no common
measure keeps only the constant''. Rational slopes close up and keep the
recipes orthogonal to them; irrational ones keep nothing but the mean.
That dichotomy, between closed orbits at rational ratios and evenly spread
ones at irrational ratios, is the same dichotomy as stations against
deserts in Part~\ref{p:cf}, seen from the torus.

\begin{checkyourself}
Why does the argument for ``linear rules see only the average'' fail for
the rule $F(I)=\int I^2$?
\answer $\int(\text{average of }I(\cdot+h))^2$ is not the average of
$\int I(\cdot+h)^2$: squaring does not move through an average. Indeed
$\int I^2$ is invariant but not a function of $\E_HI$.
\end{checkyourself}

\section{Sampling and aliasing: a strobe is a family}
\label{p:sampling}

\begin{why}
A camera taking $f$ frames a second is a family in time whose picture is a
comb of instants. The paper treats the wagon-wheel effect as a beat between
the spokes and the frames, and needs one fact: a comb has every harmonic at
equal strength.
\end{why}

\paragraph{The comb.}
Sampling a signal at times $n/f$ keeps its values at those instants only.
Viewed as a periodic function of the frame count $u=ft$, the sampler's
profile is a very narrow pulse of duty $d$ at each whole number, and
Example~\ref{e:pulse} gives its coefficients: $\hat g(n)=\sin(n\pi d)/n\pi$,
which for $d\to0$ is $d$ for every $n$ up to $|n|\sim1/d$. Divide by $d$ to
normalise, and a spike has every harmonic at strength one: the sampler is
a family whose profile carries every harmonic equally (no delta function
is needed; a spike is the limit of narrow pulses, and that limit is the
statement). So the paper's price $|ab|$ for a recipe $(a,-b)$ between a
spoke profile and the sampler charges only $|a|$, the spoke profile's
harmonic; the sampler's harmonics are free.

\paragraph{Aliasing as a beat.}
A spoke passing at rate $r$ against frames at rate $f$ has recipes
$(a,-b)$ with rate $ar-bf$. The classical alias is $a=1$: for $r$ near
$bf$ the recipe $(1,-b)$ is slow and the wheel is seen turning at $r-bf$,
backwards when $r<bf$. Nyquist's condition, that the alias rate $|r-bf|$
never fall below the true rate $r$ for any $b\ge1$, is $r<f/2$; the
wagon-wheel reversal is what happens when it is violated. With harmonics,
$a=2$: for $r$ near $f/2$ the recipe $(2,-1)$ is slow, carried by the spoke
profile's second harmonic, and a pooling eye sees a still wheel with twice
the spokes --- which a $50\%$ spoke duty extinguishes,
Example~\ref{e:pulse}. The paper's \texttt{wagonwheel.mjs} measures both.

\begin{checkyourself}
Spokes at $r=36$ per second, frames at $f=24$. What does the eye see?
\answer Recipes $(1,-1)$ at rate $36-24=12$ and $(1,-2)$ at $36-48=-12$
are equally slow; the price breaks the tie, $|ab|=1$ against $2$, so the
wheel is seen turning forwards at $12$ per second. Also $2r-3f=0$: the
recipe $(2,-3)$ is exactly slow, so if the spokes are narrow a still,
doubled wheel is superposed on the moving one; at $50\%$ duty only the
moving wheel.
\end{checkyourself}

\section{The paper's words}
\label{p:glossary}

\begin{description}\small
\item[Family] anything that comes in numbered members: ticks, clicks,
lines, rings, rays, frames.
\item[Count $\cnt$] members up to here, interpolated; Part~\ref{p:count}.
\item[Pitch] the distance between neighbouring members.
\item[Superposition] two or more families laid over each other, whether
they add (sounds, lights) or multiply (inks); the \textbf{printed pair} is
two gratings whose inks multiply.
\item[Third pattern, beat] the pattern of coincidences between two
families, on neither of them.
\item[State $\Phi$] the fractional counts of every family: a point of the
torus; Part~\ref{p:torus}.
\item[Picture $I$] what a state looks like: a function on the torus.
\item[Recipe $k$] whole numbers combining counts; a wave of the picture;
Parts~\ref{p:torus}, \ref{p:fourier}.
\item[Octave recipe] $(2,-1)$: slow when one family's rate is about twice
the other's.
\item[Order] the sum of a recipe's entries' sizes; \textbf{price} their
product; Part~\ref{p:cf}.
\item[Slow, fast] a combined count that changes slowly, or quickly,
compared with the members; Part~\ref{p:rates}.
\item[Heterodyne ratio $\het$] the rate of a combination over the average
rate of its ingredients; one over the members per band.
\item[Beat regime] $\het<\tfrac14$: more than four members per band.
\item[Priced merit $\mu$] $\het$ times the price; the visible recipe
minimises it; Part~\ref{p:cf}.
\item[Station] a recipe beyond first order that shows; it locks the pitches
at a rational ratio.
\item[Desert] a ratio with no station at any order: bounded partial
quotients.
\item[Carrier] the members themselves, the fine texture pooling removes.
\item[Slide, schedule $w$] a straight motion on the torus by whole numbers
of members.
\item[Pooling, envelope, pooled view] averaging before responding, and what
it produces: $\E[I\mid\text{slow}]$; Part~\ref{p:average}.
\item[Window $W$, response $\widehat W$] how an observer pools, and what it
does to a wave; Part~\ref{p:blur}.
\item[Front end $N$] what an observer does to each value before pooling;
Part~\ref{p:frontend}.
\item[Hard, soft] profiles taking two values, or blurred between them.
\item[Duty $d$] the fraction of a period a member covers.
\item[Duty null] a station extinguished because a profile lacks the harmonic
that carries it.
\item[Rephasing] shifting a family's count by a constant: a slide.
\item[Rung] which of the three ways a count exists: exact, winding, fold;
Part~\ref{p:implicit}.
\item[Defect, charge $q$] a point a count winds around, and by how much;
Part~\ref{p:winding}.
\item[Fold] where two members meet a point at the same index and the count
branches.
\item[Walking family] a nest whose members step in size, position and
angle.
\item[Gauge $\gamma$] a shape's own ruler.
\item[Gate] a condition a script's result must meet, written before the
result was known.
\item[Golden image] a stored render the tool must reproduce exactly on
every change.
\end{description}

\end{document}
